\begin{lemma}
\label{lemma-generator-P1}
Let $A$ be a ring. Let $X = \mathbf{P}^n_A$. Then
$$
E =
\mathcal{O}_X \oplus \mathcal{O}_X(-1) \oplus \ldots \oplus \mathcal{O}_X(-n)
$$
is a generator
(Derived Categories, Definition \ref{derived-definition-generators})
of $D_\QCoh(X)$.
\end{lemma}

\begin{proof}
Let $K \in D_\QCoh(\mathcal{O}_X)$. Assume
$\Hom(E, K[p]) = 0$ for all $p \in \mathbf{Z}$.
We have to show that $K = 0$.
By Derived Categories, Lemma
\ref{derived-lemma-right-orthogonal}
we see that $\Hom(E', K[p])$ is zero for all $E' \in \langle E \rangle$
and $p \in \mathbf{Z}$.
By Lemma \ref{lemma-construct-the-next-one}
applied with $a = -n - 1$
we see that $\mathcal{O}_X(-n - 1) \in \langle E \rangle$
because it is quasi-isomorphic to a finite complex
whose terms are finite direct sums of summands of $E$.
Repeating the argument with $a = -n - 2$ we see that
$\mathcal{O}_X(-n - 2) \in \langle E \rangle$.
Arguing by induction we find that $\mathcal{O}_X(-m) \in \langle E \rangle$
for all $m \geq 0$.
Since
$$
\Hom(\mathcal{O}_X(-m), K[p]) =
H^p(X, K \otimes_{\mathcal{O}_X}^\mathbf{L} \mathcal{O}_X(m)) =
H^p(X, K \otimes_{\mathcal{O}_X}^\mathbf{L} \mathcal{O}_X(1)^{\otimes m})
$$
we conclude that $K = 0$ by Lemma \ref{lemma-nonzero-some-cohomology}.
(This also uses that $\mathcal{O}_X(1)$ is an ample
invertible sheaf on $X$ which follows from
Properties, Lemma \ref{properties-lemma-open-in-proj-ample}.)
\end{proof}